\section*{Chapter \ref{ch:nw:buying-connectivity} --- Buying Connectivity}

\begin{solution}{exo:nw:buying-connectivity:1}
MTBF $8\,766/2 = \qty{4383}{\hour}$, MTTR \qty{4}{\hour}: availability $4\,383/4\,387 = 99.9088\%$, down $8\,766 \times 60 \times 0.000912 = 480$ minutes a
year (eight hours).
\end{solution}

\begin{solution}{exo:nw:buying-connectivity:2}
$9\,000/10 = 900$~USD per megabit a month at 10~megabits; $44\,000/200 = 220$~USD at 200.
\end{solution}

\begin{solution}{exo:nw:buying-connectivity:3}
Below 95\% of $30 \times 1\,440 = 43\,200$ minutes: more than 2\,160 minutes, 36 hours, in the month.
\end{solution}

\begin{solution}{exo:nw:buying-connectivity:4}
The pair is unavailable $u^2 + c$: with $u = 0.00091$ the independent term is $8.3 \times 10^{-7}$, negligible, and the duct's $c = 6/8\,766 = 0.00068$
is of the same order as one circuit's $u$. The second circuit removes the independent failures, which were never the problem, and keeps the common
one.
\end{solution}

\begin{solution}{exo:nw:buying-connectivity:5}
The building entry or the street near it; the \omterm{def:nw:colocation-products-and-how-they-are-sold:mmr}{meet-me room} or the carrier cage; a patch panel or a switch at either end; the power supply of the
racks; the provider itself (a software fault, a maintenance error, a bankruptcy) if both come from one; a wholesale fibre carried under two brands.
\end{solution}

\begin{solution}{exo:nw:buying-connectivity:6}
When the firm needs many venues and data sources at moderate latency, or fast access to a new venue, and does not want to build and operate
circuits to each. Not when the strategy's edge is latency to that venue: the extranet's own network and shared capacity sit in the path, and the
firm cannot control them.
\end{solution}

\begin{solution}{exo:nw:buying-connectivity:7}
Solving $u^2 + c - u^2 c = u$ with the model's circuits gives $c = 0.000911$, that is 0.67 duct events a year of twelve hours: at that rate the second
circuit in the same duct buys nothing at all.
\end{solution}

\begin{solution}{exo:nw:buying-connectivity:8}
The credit is a fraction of the month's charge for the failed service (10\% at the first threshold), credited against future bills, and it is the
sole remedy; the firm's trading loss during the outage is not covered. And 99.99\% of a month still allows four minutes of outage a month without
any credit at all.
\end{solution}

\begin{solution}{pb:nw:buying-connectivity:1}
\textbf{Part I.}
\begin{enumerate}
\item 99.9088\%, 479.6 minutes a year.
\item $12 \times (11\,000 + 600) = 139\,200$~USD a year.
\item $479.6 \times 2\,000 \approx 959\,000$~USD.
\item Nothing: four hours in a month leaves 99.44\% uptime, above the 95\% threshold.
\end{enumerate}
\textbf{Part II.}
\begin{enumerate}[resume]
\item 99.9315\%, 360.4 minutes a year.
\item 360.0 of the 360.4 minutes: half an event a year of twelve hours.
\item 278\,400~USD a year; it avoids about 238\,000~USD of expected loss against one circuit, less than its extra 139\,200 plus the risk it keeps.
\item 1\,992 minutes, 33 hours: worse than one circuit's 1\,541, because a duct event lasts twelve hours.
\end{enumerate}
\textbf{Part III.}
\begin{enumerate}[resume]
\item 99.99992\%, 0.44 minutes a year.
\item 41\,760~USD a year more than the duct; it saves 720\,000~USD of expected loss.
\item One circuit about 5.50~USD, the duct pair about 2\,560~USD, the diverse pair nothing, under the schedules of their tiers.
\item A shared element missed in the survey (building entry, power, the provider's own systems), a regional disaster, or a failure at the venue's
  end, which no circuit fixes (chapter~28).
\end{enumerate}
\textbf{Part IV.}
\begin{enumerate}[resume]
\item \emph{Named result}: two circuits in one duct are available 99.93\% of the time (360 minutes down a year, 721\,000~USD of expected loss); two
  diverse circuits 99.99992\% (0.4 minutes); the \omterm{def:nw:buying-connectivity:sla}{service credits} the duct pair earns are about 2\,560~USD a year, less than 0.4\% of its expected loss.
\item The rate and duration of common-mode events: the whole difference between the designs lives there.
\item Ask for the route maps of both circuits down to the street and the building entry, check them against each other and against the
  providers' shared infrastructure, and test by failing one circuit.
\item Latency and its stability (not only availability), time to repair, advance notice of maintenance, and the provider's diversity commitments.
\item Both active is better: traffic is spread or duplicated, failure detection is immediate, and the standby is known to work; active-standby
  needs regular fail-over tests to be trusted.
\item In its resilience section: each connection's design, availability and cost, and the loss it protects.
\item Protection against the independent failures, which were a small part of the risk; it does not protect against the duct.
\item Two things that fail separately.
\end{enumerate}
\end{solution}

\begin{solution}{iq:nw:buying-connectivity:1}
The contract's statement of performance (availability, sometimes latency and repair time), how it is measured, and what the provider owes when
it fails. A credit compensates the fee for the service not delivered, usually a fraction of a month's charge; it does not compensate the customer's
losses.
\iqlookfor{Measurement period and method; credits as a fraction of fees; sole remedy; exclusions.}
\end{solution}

\begin{solution}{iq:nw:buying-connectivity:2}
Through anything they share: a duct, a building entry, a room, a panel, power, a wholesale carrier under both, or the venue's side of the connection;
and through common causes: a regional disaster, a configuration pushed to both.
\iqlookfor{Physical and logical common modes; verification by route maps and tests.}
\end{solution}

\begin{solution}{iq:nw:buying-connectivity:3}
For breadth and speed of set-up (many venues and feeds through one connection, a new venue in weeks) when the strategy is not latency-critical at
that venue; not for the venues where the firm races.
\iqlookfor{Latency against convenience; the provider's shared network; venue agreements and fees.}
\end{solution}

\begin{solution}{iq:nw:buying-connectivity:4}
Availability is $\mathrm{MTBF}/(\mathrm{MTBF}+\mathrm{MTTR})$; with unavailability $u$ each, two independent circuits are down together $u^2$ of the time,
for example $0.001^2 = 10^{-6}$; but any common mode $c$ adds directly and usually dominates.
\iqlookfor{The formula; squares for independence; common mode as the real risk.}
\end{solution}

\begin{solution}{iq:nw:buying-connectivity:5}
For time, not bandwidth: a few megabits that arrive hundreds of microseconds or milliseconds earlier carry the messages that decide the races, and
their value exceeds the price; everything else stays on fibre.
\iqlookfor{Value of latency; small messages; fibre back-up.}
\end{solution}

\begin{solution}{iq:nw:buying-connectivity:6}
Two diverse paths from the office to each city (different providers, streets, entries, power), dual connections at each venue from two diverse
points, both active; verify route maps, test fail-over regularly, monitor latency and loss on every path, and compute the design's availability with
explicit common-mode assumptions; five nines is about five minutes a year, which only diversity gives.
\iqlookfor{Diversity end to end; active-active; testing; the arithmetic of 99.999\%.}
\end{solution}
